\section{Limites et Perspectives}
\label{sec:limites}

\subsection{Limites du Modèle}

\begin{description}
  \item[SSVI] La paramétrisation power-law peut manquer de flexibilité
    pour les options très courtes maturités (wings extrêmes).
  \item[Couverture statique] Le portefeuille est recalibré une seule fois :
    la couverture se dégrade si les grecques changent significativement.
  \item[Choc simplifié] Le décalage rigide de $\sigma_{SSVI}$ ignore la
    déformation du skew après un choc de marché.
  \item[Liquidité] Les quantités issues de l'optimisation ne tiennent pas
    compte de la profondeur réelle du carnet d'ordres.
\end{description}

\subsection{Pistes d'Amélioration}

\begin{itemize}
  \item Calibration slice-par-slice (SVI) pour meilleure précision locale.
  \item Couverture dynamique (delta-hedging continu ou discret).
  \item Calcul du VaR/CVaR du portefeuille.
  \item Scénarios de stress multi-facteurs (corrélations BTC/ETH).
\end{itemize}
